\documentclass[12pt,a4paper]{article}
\usepackage[utf8]{inputenc}
\usepackage[T1]{fontenc}
\usepackage[turkish]{babel}
\usepackage{amsmath}
\usepackage{tabularx}
\usepackage{booktabs}
\usepackage{geometry}
\geometry{a4paper, margin=1in}
\usepackage{xcolor}
\usepackage{titlesec}
\usepackage{hyperref}

\title{\textbf{İspanyolca Ders Notları - 4. Hafta}}
\author{Aula Internacional \textbf{Plus 1} (1-22. Sayfalar ve Arka Sayfa Referansları)}
\date{\today}

\begin{document}

\maketitle
\tableofcontents
\newpage

\section{Ünite 1: Nosotros (Biz) - Temel İletişim}
Bu ünitede kendimizi tanıtmayı, başkalarıyla tanışmayı, iletişim bilgilerimizi, milliyetimizi ve mesleğimizi söylemeyi; sınıftaki eşyaların isimlerini ve İspanyol alfabesini öğreniyoruz.

\subsection{Selamlaşma ve Vedalaşma (Saludos y despedidas)}
\begin{itemize}
    \item \textbf{¡Hola!}: Merhaba! (Günün her saati, en yaygın kullanım)
    \item \textbf{Buenos días}: Günaydın (Öğlene kadar)
    \item \textbf{Buenas tardes}: Tünaydın / İyi akşamüstleri (Öğleden hava kararana kadar)
    \item \textbf{Buenas noches}: İyi geceler (Hava karardıktan sonra)
    \item \textbf{¿Qué tal?}: Nasılsın?
\end{itemize}

\textbf{Vedalaşırken:}
\begin{itemize}
    \item \textbf{¡Adiós!}: Hoşça kal
    \item \textbf{¡Hasta luego!}: Sonra görüşürüz
    \item \textbf{¡Hasta pronto!}: Yakında görüşürüz
\end{itemize}

\subsection{İspanyol Alfabesi ve Okunuş (Letras y sonidos)}
İspanyolcada kelimelerin yazılışı ve okunuşu çoğunlukla aynıdır, ancak şu harflere dikkat edilmelidir:
\begin{itemize}
    \item \textbf{H (hache)}: \textcolor{red}{Okunmaz.} (\textit{Hola} $\rightarrow$ "Ola", \textit{Hotel} $\rightarrow$ "Otel")
    \item \textbf{LL (doble ele)}: Türkçe "Y" gibi okunur. (\textit{Llamo} $\rightarrow$ "Yamo")
    \item \textbf{Ñ (eñe)}: "NY" sesi verir. (\textit{España} $\rightarrow$ "Espanya")
    \item \textbf{J (jota)}: Gırtlaktan hırıltılı "H" sesiyle okunur. (\textit{Japón} $\rightarrow$ "Hapon")
    \item \textbf{G (ge)}: Yanına "e, i" alırsa hırıltılı "H" (\textit{gente} $\rightarrow$ "hente"); "a, o, u" alırsa normal "G" okunur (\textit{gato}).
    \item \textbf{C (ce)}: Yanına "e, i" alırsa Peltek S (İspanya) veya normal S (Latin Amerika) okunur (\textit{cero} $\rightarrow$ "sero"); "a, o, u" alırsa "K" okunur (\textit{casa} $\rightarrow$ "kasa").
    \item \textbf{QU (cu)}: Sadece "e" ve "i" ile kullanılır ve "K" okunur. (\textit{Qué} $\rightarrow$ "Ke", \textit{Quince} $\rightarrow$ "Kinse"). \textit{U harfi okunmaz.}
\end{itemize}

\subsection{Kişisel Bilgiler (Datos personales)}
\begin{table}[h]
\centering
\begin{tabularx}{\textwidth}{l X}
\toprule
\textbf{Soru (Sen / Siz)} & \textbf{Cevap} \\
\midrule
¿Cómo te llamas? / ¿Cómo se llama? & (Me llamo) Daniel. \\
¿De dónde eres? / ¿De dónde es? & Soy de Turquía. / Soy turco. \\
¿A qué te dedicas? / ¿En qué trabajas? & Soy profesor. / Trabajo de camarero. \\
¿Cuántos años tienes? / ¿Cuántos años tiene? & Tengo 23 años. \\
¿Tienes móvil? / ¿Tiene móvil? & Sí, es el 612 34 56 78. \\
¿Tienes correo electrónico? & Sí, es daniel@aula.com \\
\bottomrule
\end{tabularx}
\end{table}

\newpage
\section{Dilbilgisi ve Gramer (Kitabın Arka Sayfalarındaki Tüm Referanslar)}
Kitabın ilk 22 sayfasındaki konuların detaylı dilbilgisi kuralları arka sayfalarda (Más Gramática) şu şekilde açıklanır:

\subsection{1. Şahıs Zamirleri (Pronombres Personales de Sujeto)}
İspanyolcada fiil çekimleri kimin konuştuğunu belli ettiği için şahıs zamirleri genellikle cümlede kullanılmaz (sadece vurgu yapmak için kullanılır).
\begin{itemize}
    \item \textbf{Yo:} Ben
    \item \textbf{Tú:} Sen (Samimi)
    \item \textbf{Él / Ella / Usted:} O (Erkek) / O (Kadın) / Siz (Tekil, Resmi hitap)
    \item \textbf{Nosotros / Nosotras:} Biz (Erkekler veya karışık) / Biz (Sadece kadınlar)
    \item \textbf{Vosotros / Vosotras:} Sizler (Samimi çoğul - Sadece İspanya'da kullanılır)
    \item \textbf{Ellos / Ellas / Ustedes:} Onlar (Erkekler) / Onlar (Kadınlar) / Sizler (Çoğul resmi. Ancak Latin Amerika'da herkes için "Sizler" anlamında kullanılır.)
\end{itemize}

\subsection{2. Düzenli Fiiller: -AR, -ER, -IR (El Presente de Indicativo)}
İspanyolcada fiiller 3 ana gruba ayrılır. Şimdiki/geniş zaman (Presente) çekimi, kelimenin sonundaki -ar, -er, -ir atılıp köke şu eklerin getirilmesiyle yapılır:

\begin{table}[h]
\centering
\begin{tabularx}{\textwidth}{l X X X}
\toprule
\textbf{Kişi Zamiri} & \textbf{-AR (Trabajar - Çalışmak)} & \textbf{-ER (Leer - Okumak)} & \textbf{-IR (Escribir - Yazmak)} \\
\midrule
\textbf{Yo} & trabaj\textbf{-o} & le\textbf{-o} & escrib\textbf{-o} \\
\textbf{Tú} & trabaj\textbf{-as} & le\textbf{-es} & escrib\textbf{-es} \\
\textbf{Él/Ella/Usted} & trabaj\textbf{-a} & le\textbf{-e} & escrib\textbf{-e} \\
\textbf{Nosotros/as} & trabaj\textbf{-amos} & le\textbf{-emos} & escrib\textbf{-imos} \\
\textbf{Vosotros/as} & trabaj\textbf{-áis} & le\textbf{-éis} & escrib\textbf{-ís} \\
\textbf{Ellos/Ellas/Ustedes} & trabaj\textbf{-an} & le\textbf{-en} & escrib\textbf{-en} \\
\bottomrule
\end{tabularx}
\end{table}

\subsection{3. Temel Düzensiz Fiiller: Ser, Tener, Llamarse}
\begin{table}[h]
\centering
\begin{tabularx}{\textwidth}{l X X X}
\toprule
\textbf{Kişi Zamiri} & \textbf{SER (Olmak)} & \textbf{TENER (Sahip olmak)} & \textbf{LLAMARSE (Adlandırılmak)} \\
\midrule
\textbf{Yo} & soy & tengo & me llamo \\
\textbf{Tú} & eres & tienes & te llamas \\
\textbf{Él / Ella / Usted} & es & tiene & se llama \\
\textbf{Nosotros/as} & somos & tenemos & nos llamamos \\
\textbf{Vosotros/as} & sois & tenéis & os llamáis \\
\textbf{Ellos/Ellas/Ustedes} & son & tienen & se llaman \\
\bottomrule
\end{tabularx}
\end{table}

\subsection{4. Eril ve Dişil Kuralları (El Género: Profesiones y Nacionalidades)}
İspanyolcada her ismin ve sıfatın cinsiyeti vardır.
\begin{itemize}
    \item \textbf{Eril (Masculino):} Genellikle \textbf{-o} ile biter (ör. \textit{italiano, cocinero}). Sessiz harfle bitenler de genelde erildir (ör. \textit{alemán, profesor}).
    \item \textbf{Dişil (Femenino):} Genellikle \textbf{-a} ile biter (ör. \textit{italiana, cocinera}). Sessiz harfle biten erilleri dişil yapmak için sonuna \textbf{-a} eklenir (ör. \textit{alemana, profesora}).
    \item \textbf{Değişmeyenler (Ortak):} Sonu \textbf{-e}, \textbf{-ista} veya \textbf{-í} ile biten kelimeler hem kadın hem erkek için \textbf{aynıdır} (ör. \textit{estudiante, periodista, marroquí}).
\end{itemize}

\subsection{5. Tanımlıklar (Artículos: El, La, Los, Las / Un, Una)}
Sınıftaki eşyalar (Las cosas de la clase) konusuyla birlikte isimlerin önüne gelen tanımlıklar öğrenilir.
\begin{itemize}
    \item \textbf{Belirli (Definidos):} Bilinen bir nesneden bahsederken kullanılır (İngilizcedeki "the").
    \begin{itemize}
        \item Eril Tekil: \textbf{el} (el bolígrafo) \textbar{} Eril Çoğul: \textbf{los} (los libros)
        \item Dişil Tekil: \textbf{la} (la mesa) \textbar{} Dişil Çoğul: \textbf{las} (las sillas)
    \end{itemize}
    \item \textbf{Belirsiz (Indefinidos):} Herhangi bir nesneden bahsederken kullanılır (İngilizcedeki "a/an").
    \begin{itemize}
        \item Eril Tekil: \textbf{un} (un cuaderno) \textbar{} Eril Çoğul: \textbf{unos} (unos cuadernos)
        \item Dişil Tekil: \textbf{una} (una hoja) \textbar{} Dişil Çoğul: \textbf{unas} (unas hojas)
    \end{itemize}
\end{itemize}

\subsection{6. Soru Kelimeleri (Los Interrogativos)}
Soru sorarken (¿?) başa gelen soru kelimeleri her zaman \textbf{aksanlı (üzeri çizgili)} yazılır.
\begin{itemize}
    \item \textbf{¿Qué?:} Ne? (Örn: \textit{¿Qué significa?})
    \item \textbf{¿Cómo?:} Nasıl? (Örn: \textit{¿Cómo te llamas?, ¿Cómo se dice?})
    \item \textbf{¿De dónde?:} Nereden? (Örn: \textit{¿De dónde eres?})
    \item \textbf{¿Cuántos/Cuántas?:} Kaç? / Ne kadar? (Örn: \textit{¿Cuántos años tienes?})
    \item \textbf{¿Cuál?:} Hangisi? / Nedir? (Örn: \textit{¿Cuál es tu correo electrónico?})
    \item \textbf{¿Quién?:} Kim? (Örn: \textit{¿Quién es?})
\end{itemize}
\textit{Not:} "Ustedes son estudiantes?" cümlesinde olduğu gibi, soru kelimesi olmadan da sadece vurgu ve noktalama ile soru cümlesi kurulabilir.

\subsection{7. Sayılar (Los números del 0 al 100)}
\begin{itemize}
    \item \textbf{0-10:} cero (0), uno (1), dos (2), tres (3), cuatro (4), cinco (5), seis (6), siete (7), ocho (8), nueve (9), diez (10).
    \item \textbf{11-20:} once (11), doce (12), trece (13), catorce (14), quince (15), dieciséis (16), diecisiete (17), dieciocho (18), diecinueve (19), veinte (20).
    \item \textbf{21-29 (Bitişik Yazılır):} veintiuno (21), veintidós (22), veintitrés (23), ... veintinueve (29).
    \item \textbf{Onluklar:} treinta (30), cuarenta (40), cincuenta (50), sesenta (60), setenta (70), ochenta (80), noventa (90).
    \item \textbf{31'den İtibaren (Ayrı Yazılır):} Onluk ve birlik arasına "y" (ve) harfi konur. \textit{Örn: treinta y uno (31), cuarenta y cinco (45), ochenta y dos (82).}
    \item \textbf{100:} cien.
\end{itemize}

\newpage
\newpage
\section{Pratik ve Alıştırmalar (Ejercicios)}
Madem notlardan ilerliyoruz, kitabı açmana gerek kalmadan doğrudan burada çözebileceğin, ilk ünitenin temel kazanımlarını ölçen alıştırmaları senin için hazırladım. (Cevap anahtarı bir sonraki sayfadadır).

\subsection{Alıştırma 1: Sınıf Eşyaları ve Tanımlıklar (El / La)}
Aşağıdaki kelimelerin cinsiyetini (Eril/Dişil) tahmin et ve boşluklara uygun tanımlığı (\textbf{el} veya \textbf{la}) yerleştir. Kuralı hatırla: Genelde -o ile bitenler \textit{el}, -a ile bitenler \textit{la} alır.
\begin{enumerate}
    \item \_\_\_\_\_\_\_ bolígrafo (Tükenmez kalem)
    \item \_\_\_\_\_\_\_ mesa (Masa)
    \item \_\_\_\_\_\_\_ libro (Kitap)
    \item \_\_\_\_\_\_\_ silla (Sandalye)
    \item \_\_\_\_\_\_\_ cuaderno (Defter)
    \item \_\_\_\_\_\_\_ pizarra (Yazı tahtası)
\end{enumerate}

\subsection{Alıştırma 2: Fiil Çekimleri (Ser, Tener, Llamarse)}
Boşlukları parantez içindeki fiilin doğru şahıs çekimiyle doldur.
\begin{enumerate}
    \item ¡Hola! Yo \_\_\_\_\_\_\_\_\_\_\_\_\_\_ (llamarse) Alejandro.
    \item ¿De dónde \_\_\_\_\_\_\_\_\_\_\_\_\_\_ (ser) tú?
    \item Nosotros \_\_\_\_\_\_\_\_\_\_\_\_\_\_ (tener) 25 años.
    \item ¿A qué \_\_\_\_\_\_\_\_\_\_\_\_\_\_ (dedicarse) usted?
    \item Ellos \_\_\_\_\_\_\_\_\_\_\_\_\_\_ (ser) de México.
    \item ¿(Tú) \_\_\_\_\_\_\_\_\_\_\_\_\_\_ (tener) móvil?
\end{enumerate}

\subsection{Alıştırma 3: Soru Kelimeleri (Los Interrogativos)}
Aşağıdaki diyalogları uygun soru kelimeleriyle (\textbf{Cómo, De dónde, Cuántos, Qué, Cuál}) doldur.
\begin{enumerate}
    \item — ¿\_\_\_\_\_\_\_\_\_\_\_\_ te llamas? \\ — Me llamo Carmen.
    \item — ¿\_\_\_\_\_\_\_\_\_\_\_\_ eres? \\ — Soy de Colombia.
    \item — ¿\_\_\_\_\_\_\_\_\_\_\_\_ años tienes? \\ — Tengo 21 años.
    \item — ¿\_\_\_\_\_\_\_\_\_\_\_\_ es tu correo electrónico? \\ — Es carmen@aula.com.
    \item — ¿\_\_\_\_\_\_\_\_\_\_\_\_ significa "ordenador"? \\ — "Bilgisayar" demektir.
\end{enumerate}

\subsection{Alıştırma 4: ¿Ustedes son estudiantes?}
Aşağıdaki durumlarda hangi hitap şeklini (\textbf{tú, usted, vosotros, ustedes}) kullanman gerektiğini yaz.
\begin{enumerate}
    \item Yeni tanıştığın, senden yaşça büyük birine "Siz" diyeceksin: \_\_\_\_\_\_\_\_
    \item Yakın arkadaşınla konuşuyorsun: \_\_\_\_\_\_\_\_
    \item (İspanya'da) Sınıftaki arkadaşlarına (çoğul) sesleniyorsun: \_\_\_\_\_\_\_\_
    \item (Latin Amerika'da) Sınıftaki arkadaşlarına (çoğul) sesleniyorsun: \_\_\_\_\_\_\_\_
\end{enumerate}

\newpage
\section*{Cevap Anahtarı (Clave de respuestas)}
\textbf{Alıştırma 1:} 1. El, 2. La, 3. El, 4. La, 5. El, 6. La \\
\textbf{Alıştırma 2:} 1. me llamo, 2. eres, 3. tenemos, 4. se dedica, 5. son, 6. tienes \\
\textbf{Alıştırma 3:} 1. Cómo, 2. De dónde, 3. Cuántos, 4. Cuál, 5. Qué \\
\textbf{Alıştırma 4:} 1. usted, 2. tú, 3. vosotros, 4. ustedes

\end{document}
